% 11-711 Fall 2026; Homework 2 report template.
%
% This is the ACL template, which the handout requires. Keep the style files as
% they are: do not change margins, font size, or spacing to fit more in.
%
% HARD LIMIT: 7 pages of content. References do not count toward the limit.
%
% The four graded sections below are worth 10 points each (40 of 100). Each
% one lists the questions the handout asks you to answer; answer all of them.
% If something does not apply to your system, say so in a sentence rather than
% leaving the section empty, so we can tell the difference.
%
% Cite every model, library, dataset, paper, and blog post you used. Per the
% syllabus, every word of this report must be your own: you may use AI to help
% write code, but not to write any part of this report.

\pdfoutput=1

\documentclass[11pt]{article}

\usepackage{acl}

\usepackage{times}
\usepackage{latexsym}
\usepackage{amsmath}
\usepackage{array}
\usepackage{tabularx}
\usepackage{booktabs}
\usepackage{multirow}
\usepackage{graphicx}
\usepackage[T1]{fontenc}
\usepackage[utf8]{inputenc}
\usepackage{microtype}
\usepackage{inconsolata}

\title{11-711 Homework 2: Retrieval-Augmented Generation for Pittsburgh and CMU}

% Individual submission. If your team is approved to work jointly, list every
% author and mark equal contribution as needed.
\author{
  Your Name \\
  \texttt{yourandrewid@andrew.cmu.edu} \\[4pt]
  {\small Due: October 23, 2026, 11:59 PM ET}
}

\begin{document}
\maketitle

\begin{abstract}
% Around 150 words. What system did you build, what were the main design
% decisions, and what did the evaluation show? State your headline number.
\end{abstract}

%======================================================================
\section{Introduction}
% Replace this sentence; it is here to show the citation commands and to give
% the bibliography something to render. \citep for parenthetical, \citet for
% textual: as \citet{lewis2021rag} show, ...
We build a retrieval-augmented generation system \citep{lewis2021rag} that
answers factual questions about Pittsburgh and Carnegie Mellon University.

% Half a page at most. What the task is, what your final system does end to
% end, and a roadmap of the report. A pipeline figure here is usually worth
% more than a paragraph of prose.

% \begin{figure}[t]
%   \centering
%   \includegraphics[width=\columnwidth]{assets/pipeline.png}
%   \caption{End-to-end architecture of our system.}
%   \label{fig:pipeline}
% \end{figure}

%======================================================================
\section{Data Creation}
\label{sec:data}
% 10 points. Answer all four questions below.

\subsection{Knowledge resource}
% How did you compile your knowledge resource, and how did you decide which
% documents to include? Give the source list and the scoping rule you applied.
% A table of source, number of documents, and number of tokens makes the scale
% concrete and costs you three lines.

% \begin{table}[t]
%   \centering
%   \begin{tabular}{lrr}
%     \toprule
%     Source & Documents & Tokens \\
%     \midrule
%     & & \\
%     \midrule
%     Total & & \\
%     \bottomrule
%   \end{tabular}
%   \caption{Composition of our knowledge resource.}
%   \label{tab:corpus}
% \end{table}

\subsection{Extraction and cleaning}
% How did you extract the raw data, and with what tools? Cover crawling, HTML
% and PDF handling, boilerplate removal, and deduplication. Report what you
% discarded and why, not only what you kept.

\subsection{Annotated training data}
% What did you annotate, what kind, and how much? Describe your guidelines and
% your annotator agreement if you measured it. This step is optional in the
% assignment: if you annotated nothing, say so here and explain the choice.

\subsection{Other data}
% For training data you did not annotate yourself: what extra data did you use,
% and in what way? Note licenses, and confirm it falls within the model and
% data policy in the handout.

%======================================================================
\section{Model Details}
\label{sec:model}
% 10 points. Describe your system, and explain at least TWO variations you
% tried (more is welcome) along with your justification for trying each. A
% variation can be a different embedding model, retriever, re-ranker, chunking
% scheme, reader, prompt, or training strategy.

\subsection{Chunking}
% Your splitting logic, chunk size, overlap, and how you handled boundaries.
% Say how you treated the edge cases: very short documents, very long
% paragraphs, tables and other special formatting.

\subsection{Dense retrieval}
% Embedding model and why you chose it, index type, and how you build and query
% it. Report the number of vectors and the retrieval depth $k$.

\subsection{Sparse retrieval}
% BM25 implementation, tokenization and stemming choices, and any parameter
% tuning.

\subsection{Hybrid retrieval}
% The combination strategy is a required component. State the fusion method
% precisely enough that a reader could reimplement it: give the formula, the
% weights, and how you chose them.

\subsection{Reader}
% The generation model, prompt format, decoding settings, and context budget.
% Confirm the model satisfies the handout's policy (open-weight, on HuggingFace,
% within the size limit).

\subsection{Baselines}
% What you compare against. At minimum a closed-book variant of your reader,
% since Section~\ref{sec:analysis} asks you to quantify what retrieval buys.

%======================================================================
\section{Results}
\label{sec:results}
% 10 points. Report the raw numbers for every system variant you tried, measured
% on public_evaluation_data.json and on any development or test data you built
% yourself. Add the leaderboard score for the systems whose predictions you
% submitted there; with 10 submissions you will not have one for every variant,
% and that is expected. State which system produced each of the
% system_output_{1,2,3}.json files you submitted.

% \begin{table}[t]
%   \centering
%   \begin{tabular}{lrrrrr}
%     \toprule
%     & \multicolumn{4}{c}{Public eval (157 q)} & Leaderboard \\
%     \cmidrule(lr){2-5}
%     System & Recall & F1 & ROUGE-L & Judge & Total \\
%     \midrule
%     Closed-book      & & & & & -- \\
%     Dense only       & & & & & -- \\
%     Sparse only      & & & & & -- \\
%     Hybrid (final)   & & & & & \\
%     \bottomrule
%   \end{tabular}
%   \caption{Results on the public evaluation set, with leaderboard scores for the
%            systems we submitted.}
%   \label{tab:results}
% \end{table}

% Say which set each number comes from and how many questions it covers.
% Leaderboard scores and numbers from your own held-out set measure different
% questions, so keep them in separate columns and never average the two.

%======================================================================
\section{Analysis}
\label{sec:analysis}
% 10 points. This section is graded on insight, not on the size of the table.
% All five items below are required.

\subsection{Performance by question type}
% Break the results down more finely than one holistic number. Define your
% question categories, say how you assigned them, and report per-category
% scores. Where does the system fail, and why?

\subsection{Comparing your variations}
% Report results for at least two of the variations from Section~\ref{sec:model}
% and say what you conclude. If two systems differ by less than run-to-run
% variation, say that rather than declaring a winner.

\subsection{Retrieval versus closed-book}
% Quantify what retrieval actually buys over using the same reader with no
% context. Which questions improve, and which get worse?

\subsection{Dense, sparse, and hybrid}
% Compare the three retrieval modes directly, and compare your fusion
% strategies. Which works best for which kind of question? A question type
% where sparse beats dense, or the reverse, is the interesting part.

\subsection{Example outputs}
% Show outputs from at least two of your systems on the same questions. Choose
% examples that illustrate the quantitative differences above rather than
% cherry-picking successes.

%======================================================================
\section{Conclusion}
% What you built, what you learned, and what you would do next with more time.

%======================================================================
\section*{Limitations}
% Optional but encouraged: where your evaluation is thin, what your knowledge
% resource does not cover, and which conclusions your evidence does not support.

%======================================================================
\section*{Use of AI (required, ungraded)}
% Which parts of this assignment did you use AI for, and how? If you used none,
% say so. This question is required by the syllabus: using AI without making an
% honest attempt at answering it is itself an academic integrity violation.
% Remember that AI may not be used to write any part of this report.

%======================================================================
\bibliography{custom}

\end{document}
